\chapter*{Résumé}
\addcontentsline{toc}{chapter}{Résumé}

Ce rapport évalue la performance du régime fiscal minier applicable au secteur aurifère sénégalais à l'aide d'une adaptation du modèle \emph{Fiscal Analysis of Resource Industries} (FARI) du Fonds Monétaire International, repris dans l'outil de modélisation financière d'une mine d'or diffusé par le \emph{Natural Resource Governance Institute} (NRGI). Le projet retenu pour l'application est le complexe aurifère Sabodala-Massawa, exploité par Sabodala Gold Operations (SGO), filiale d'Endeavour Mining (ex-Teranga Gold), dont les données techniques et économiques sont tirées de l'étude de préfaisabilité (PFS) publiée en août~2020.

Les données du projet (réserves, calendrier de production, dépenses d'investissement, coûts d'exploitation, prix de l'or retenu) ont été extraites du PFS et insérées dans le modèle, dont l'architecture --- assiette des redevances, impôt sur les bénéfices, participation de l'État, retenues à la source, droits de douane --- a par ailleurs été complétée par un régime fiscal ``Sénégal'' reproduisant les dispositions effectivement applicables à SGO~: redevance de 5\,\% sur la valeur des expéditions d'or, impôt sur les sociétés de 25\,\%, participation gratuite de l'État de 10\,\% et exonération de droits de douane au titre du statut d'entreprise franche d'exportation.

Les résultats indiquent que le régime fiscal applicable à SGO génère, sur la durée de vie résiduelle du projet (2020-2036), des recettes publiques représentant environ 44\,\% de la valeur du flux de trésorerie avant impôt du projet (Taux Effectif Moyen d'Imposition, \TEMI), soit un niveau proche de celui observé en République démocratique du Congo avant l'amendement de 2018 (44,5\,\%), mais nettement inférieur à celui du code minier congolais amendé (57,6\,\%) et à la moyenne des régimes comparés (45-55\,\%), à l'exception notable du Pérou et de la Tanzanie. La structure des recettes publiques sénégalaises repose principalement sur l'impôt sur les bénéfices (49\,\%), la redevance (22\,\%), la participation gratuite de l'État (16\,\%) et la retenue à la source sur les dividendes (14\,\%), l'absence de droits de douane --- liée au régime d'entreprise franche --- constituant la principale source d'écart avec les régimes comparateurs.

Une limite méthodologique importante est mise en évidence~: le profil de trésorerie de Sabodala-Massawa, projet d'extension d'une mine déjà en production, ne comporte pas de flux négatif initial, ce qui rend les indicateurs de taux de rentabilité interne (\TRI) du modèle --- conçus pour des projets ``greenfield'' --- non interprétables, sans pour autant affecter le calcul de la valeur actuelle nette ni du \TEMI. Le rapport se conclut par une série de recommandations destinées à l'ITIE Sénégal pour le suivi de la performance fiscale du secteur aurifère et la divulgation des paramètres économiques sous-jacents aux conventions minières.

\bigskip
noindent\textbf{Mots-clés~:} fiscalité minière, modèle FARI, Sénégal, or, Sabodala-Massawa, taux effectif moyen d'imposition, ITIE.
